\section{Introduction}\label{sec:intro}

Korean company law lets a listed company divide itself in two ways, and the difference is a difference in who decides what happens next. In an equity spin-off the new company's shares go to the parent's shareholders, as in a US spin-off. In a \captive{} (\emph{mulchŏk bunhal}) the parent keeps all of them, so the business the shareholders used to own directly they now own through the parent, and whether it stays there is decided later by the parent's board, which the controlling owner appoints, without a vote of the parent's shareholders. The Korea Exchange counted \KLKrxCaptive{} such divisions among \KLKrxTotal{} between 1998 and 2022 \cite{kang2025}; among the \NClean{} divisions of 2001--2026 whose later reports we could read, the parent went on to sell a majority of the subsidiary, list it or dilute its stake in \CleanDisposedShare{} percent of cases. The form is a two-step transaction whose second step is discretionary. The controversy of 2020--2022---LG Chem's battery business carved out in September 2020 and listed in January 2022 while LG Chem's own stock lost more than \KLLGDrop{} percent \cite{kang2025, park2022}---led the Financial Services Commission in September 2022 to a package of three measures: a mandatory statement of the division's purpose and of the measures taken to protect shareholders, an appraisal right for dissenting shareholders, and an exchange review of any listing of the subsidiary within five years \cite{fsc2022a}. The disclosure rule was written for this one conflict, and it arrived together with a right and a review.

The literature on controlling owners has measured expropriation at the transaction that carries it out. The intention is inferred from the market's reaction to that transaction: to the merger between chaebol affiliates in which the acquiring firm's minority lost while the controlling family gained through other firms of the group \cite{bae2002}, to the private securities offering priced in the controlling shareholder's favour \cite{baek2006}, to the loan to the weak affiliate \cite{jiang2010, johnson2000, bertrand2002}. Those measurements are taken when the transaction is announced; ours is taken before it, from a document filed when the transaction was still a possibility. We build an \emph{ex ante} measure for one class of transactions: what a reader of the division filing could infer, on the day it was filed, about whether the subsidiary would be kept. The division filing is a long structured document the parent must file to obtain the shareholders' vote, and a controller who means to sell or list the subsidiary has ordinary commercial reasons to say things in it that a controller who means to keep it does not: Doosan's 2021 filing said the parent would sell its entire stake in the new company to an affiliate, CJ ENM's 2020 filing gave the attraction of outside investment, and a share sale where necessary, among the division's purposes, and DB HiTek's 2023 filing wrote a shareholder vote before any listing into both companies' charters. We have two large language models read the filings of \NScored{} \captive{}s announced between 2001 and 2026, each as of its date, and forecast whether the parent will still own all of the new company three to five years later; we then trace what the parent actually did through its subsequent annual reports. We argue that the forecast measures the disposal risk the filing disclosed, and we test it the way one would test any measure: against the realised outcome, against the hand-coded purpose categories that \citeA{kang2025} used for the same population, against its own extracted facts, across subgroups of firms, and with every name and date removed from the text. The controller's private benefit and the investors' reading are not observed, and we do not ask why the parent chose the captive form over the equity form, which \citeA{kang2025a} do. We take the form as given and ask three questions. Did the division filing disclose the disposal before it happened? Did the market price what the filing disclosed? What did the 2022 reform change?

The filing did disclose it. Across the \NGateOne{} labelled divisions the forecast orders disposed from retained subsidiaries with an AUC of \AUCPooledLLM{}, against \AUCPooledFin{} for the filing's financial covariates and \NestedMTwo{} for the purpose categories and mention flags a hand coder would record; it doubles the hazard of disposal per standard deviation, and it is unchanged, in accuracy and in level, when the filings are anonymised. The market did not price it at the announcement, as far as returns can show: the three-day return rises with the forecast while the market-adjusted return over the following year falls with it, by \RankSpecThreePctl{} percentiles of the year's return ranking between the two ends of the forecast, against market, size-and-book-to-market portfolio and propensity-matched benchmarks though not against a single matched firm, and the loss is realised in the subsidiaries that changed hands within a year. Smaller controlling blocks go with stronger forecasts and more disposals: \LowBlockDisp{} percent of subsidiaries disposed of in the lowest tercile of the block against \HighBlockDisp{} in the highest, and a forecast \ScoreOnBlockTenPPSDAbs{} standard deviations lower per ten percentage points of the block. The reform's three measures left separable traces. Against equity spin-offs, which the reform did not touch, the announcement reaction to a \captive{} rose by \DiDCxP{} percentage points, while the mean forecast moved from \ScorePre{} to \ScorePost{} on \ProtPostN{} post-reform filings; the appraisal right's moneyness is priced with the sign a put implies, imprecisely in the three-day window; and the forecast is lower, and early disposal rarer, in the post-reform filings that commit the parent to keeping the subsidiary than in the \ProtNone{} that state no measure.

The contribution is the measure and what it lets us see. For the principal--principal literature \cite{young2008, claessens2002, masulis2011} it separates, in this class of transactions, what the filing disclosed about the disposal from the disposal that later revealed it, and shows that the disclosure was in the filing the shareholders voted on: Doosan's and CJ ENM's filings described the sale before the vote. For the literature on disclosure as a governance instrument, in this journal and elsewhere \cite{huang2024, liu2026, shim2024}, it distinguishes a disclosure that was \emph{available} from one that was \emph{used}: the information that forecast the minority's loss was in the filing before anyone was required to put it there, and within the horizon we observe the 2022 mandate changed what filings say and how they are received, with no detectable change in what they forecast. For regulators the result separates what each instrument changed: the statements and their reception, the price of dissent, and the content of the commitments parents now write; it does not show that controllers' conduct changed, and we say where the data cannot yet speak. Language-model readings of disclosures are entering empirical accounting and finance \cite{frankel2022, siano2025}, and Section~\ref{sec:validity} sets out the validity and bias tests we think such a measure owes its readers, and reports where ours is weaker: for small parents and for the oldest filings.
